% =============================================================================
%  LECTURE NOTES --- SECTION 5
%  "How Close Are We to AGI? A Reality Check for Decision-Makers"
%  ETH Zurich Executive Education --- 90-minute lecture
% =============================================================================
\documentclass[11pt,a4paper]{article}

\usepackage[T1]{fontenc}
\usepackage{lmodern}
\usepackage[margin=2.5cm]{geometry}
\usepackage{amsmath,amssymb,amsthm}
\usepackage{booktabs}
\usepackage{enumitem}
\usepackage{hyperref}
\usepackage{xcolor}
\usepackage{parskip}
\usepackage{tcolorbox}
\usepackage{microtype}
\usepackage[round]{natbib}

\definecolor{accent}{HTML}{1B1F3B}
\definecolor{accent2}{HTML}{00718B}

\hypersetup{
  colorlinks=true,
  linkcolor=accent,
  citecolor=accent2,
  urlcolor=accent2,
}

\tcbuselibrary{skins, breakable}
\newtcolorbox{keyinsight}[1][]{%
  enhanced, breakable,
  colback=accent!5, colframe=accent!70!black,
  boxrule=0.4pt, arc=2pt,
  left=8pt, right=8pt, top=6pt, bottom=6pt,
  title={\bfseries Key insight}, fonttitle=\bfseries,
  coltitle=white, colbacktitle=accent,
  #1
}

\title{\textbf{Lecture Notes: How Close Are We to AGI?}\\[0.3em]
  \large A Reality Check for Decision-Makers}
\author{ETH Zurich Executive Education}
\date{2026}

\begin{document}
\maketitle
\tableofcontents
\newpage

% =============================================================================
\section*{Preface: Who these notes are for}
\addcontentsline{toc}{section}{Preface}
% =============================================================================

These notes accompany a 90-minute lecture for senior managers at ETH
Zurich. The audience is analytically sharp, professionally accomplished,
and --- by design --- \emph{not} specialised in machine learning. You
will not be asked to read a gradient or inspect a loss curve. You
\emph{will} be asked to read a research paper the way a good board
member reads an audit: to see the shape of the argument, the quality of
the evidence, and the size of the claim that is actually being made.

The goal is not to tell you when Artificial General Intelligence (AGI)
will arrive. Nobody credible knows. The goal is to give you a
reference-backed briefing on where the field stands, grounded in six
foundational papers, so that you can judge future claims for yourself
and ask the right questions of the people who build, buy, or regulate
these systems on your behalf.

If you remember only one thing from the whole session, let it be this:
\emph{the question is not whether AGI is coming; it is whether your
organisation has the evaluation capability to tell, month by month,
workflow by workflow, what these systems can and cannot do for you.}

% =============================================================================
\section{What Do We Mean by AGI?}\label{sec:act1}
% =============================================================================

\subsection{Why the definition matters --- commercially}

The phrase ``artificial general intelligence'' appears in press
releases, investor decks, and legal contracts. It is also, as recently
as five years ago, a phrase most serious researchers avoided in public.
Two things changed. First, GPT-4 was surprising enough that
Microsoft Research felt compelled to publish a 155-page technical
report titled \emph{Sparks of Artificial General Intelligence}
\citep{bubeck2023sparks}. Second, Google DeepMind --- whose explicit
corporate mission is to build AGI --- decided the community needed a
shared vocabulary to argue over. The result was the paper that frames
this whole section: \emph{Levels of AGI for Operationalizing Progress
on the Path to AGI} \citep{morris2024levels}.

For a manager, the definition matters for a concrete reason: procurement
contracts, insurance policies, acceptable-use clauses, and regulatory
thresholds are starting to invoke ``AGI'' or ``frontier AI'' as a
trigger. If the term is fuzzy, the obligation is fuzzy. DeepMind's
framework is useful precisely because it trades romance for
measurability.

\subsection{The Morris et al.\ framework}

\citet{morris2024levels} observe that prior definitions of AGI ---
``human-level intelligence,'' ``passing the Turing test,''
``economically valuable work'' --- are either unfalsifiable, outdated,
or both. They propose a taxonomy organised along two axes:

\begin{itemize}[leftmargin=1.4em]
  \item \textbf{Performance:} how well the system performs relative to
        humans, ranging from \emph{Level 0} (no AI) through
        \emph{Emerging} (roughly an unskilled human), \emph{Competent}
        (50th percentile), \emph{Expert} (90th), \emph{Virtuoso}
        (99th), up to \emph{Superhuman} (outperforms all humans).
  \item \textbf{Generality:} whether the system is \emph{narrow}
        (bounded to a specific task or domain) or \emph{general}
        (capable across a wide range of non-cognitive and cognitive
        tasks).
\end{itemize}

That gives a \(5\times 2\) matrix. Historical milestones populate it.
Deep Blue and AlphaFold sit in the narrow column, as virtuoso or
superhuman on their particular task, but zero on generality. Siri and
early chatbots were general but emerging at best. The authors argue
that today's frontier large language models (they discuss GPT-4,
Gemini, Claude) sit in the general column but only at the
\emph{Emerging} level: they do many things that an unskilled adult
could do, across a wide range of domains, but they do not reliably
reach even competent human performance across the board. That is the
paper's headline diagnosis, and it is worth stating clearly: \emph{by
DeepMind's own framework, there is nothing at the `Competent General
AGI' level yet; we are still at `Emerging General AGI'\/}.

The framework also articulates six design principles that are worth
reading even if you never look at the matrix again:

\begin{enumerate}[leftmargin=1.4em]
  \item Focus on \emph{capabilities}, not on \emph{mechanisms}. What
        the system can do is measurable; whether it does so via
        ``real'' understanding is not.
  \item Focus on \emph{generality and performance} separately.
        Collapsing them hides where progress actually is.
  \item Focus on \emph{cognitive} (and optionally metacognitive) tasks,
        not physical embodiment. An AGI need not have a body.
  \item Focus on \emph{potential}, not deployment. What the model could
        do, not what it is currently permitted to do.
  \item Use \emph{ecological validity}: measure on tasks that matter in
        the real world, not just academic benchmarks.
  \item Treat AGI as a \emph{path}, not a single destination. There
        will be levels, and each level has its own governance question.
\end{enumerate}

\subsection{The counterweight: Chollet's critique}

No briefing on this topic is honest without the other side. Fran\c{c}ois
Chollet's \emph{On the Measure of Intelligence} \citep{chollet2019measure}
argues that benchmark performance --- the raw material of the Morris
taxonomy --- is a misleading indicator. Real intelligence, Chollet
argues, is \emph{the efficiency with which a system acquires new skills
in a previously unseen domain, using limited prior knowledge and
limited experience}. By that measure, a system that scores highly on
tasks statistically similar to its training data has demonstrated
\emph{memorisation plus interpolation}, not intelligence. Chollet's
ARC-AGI benchmark, designed to resist this kind of interpolation, was
only substantially cracked by frontier models in late 2024 --- and even
then at great computational expense.

The responsible reading is not to pick a side. It is to hold both
frames at once:

\begin{keyinsight}
Morris et al.\ tell you \emph{how to locate today's systems on a shared
axis}. Chollet tells you \emph{how to test whether the location is
real, or whether it reflects training-data overlap}. In commercial
settings, you need both: the location to decide scope, and the
resistance-to-interpolation test to decide trust.
\end{keyinsight}

\subsection{Managerial implications of Act 1}

Three takeaways you can use on Monday morning:

\begin{enumerate}[leftmargin=1.4em]
  \item When a vendor says ``our system is approaching AGI,'' the
        correct first question is: \emph{on which axis?} Performance
        or generality? Against what reference distribution? Without an
        answer, the claim is marketing.
  \item Today's frontier models are, by the most careful published
        framework, at the \emph{Emerging General AGI} level. That is
        both more and less than the public conversation suggests. More:
        it is genuinely a new category of artefact. Less: it is not
        yet competent across the board, and the gap between ``sometimes
        surprising'' and ``reliably useful'' is where most enterprise
        pilot projects still fail.
  \item The path is continuous, not discrete. You will not receive a
        press release announcing AGI. You will, if you are paying
        attention, notice a series of small capability jumps in
        workflows you already run.
\end{enumerate}

% =============================================================================
\section{Sparks and Specialist Breakthroughs}\label{sec:act2}
% =============================================================================

\subsection{The paper that broke the taboo: \emph{Sparks}}

In March 2023, a team of senior Microsoft Research scientists
published \emph{Sparks of Artificial General Intelligence: Early
experiments with GPT-4} \citep{bubeck2023sparks}. The paper is 155
pages of qualitative probes: draw a unicorn in TikZ code without seeing
an image, solve a novel maths olympiad problem, debug human-written
code by inferring intent, write and then critique a stanza in the
voice of a specified poet. The authors' conclusion, which caused the
paper's shockwave, was that GPT-4 shows ``sparks of artificial general
intelligence'' and that they ``could reasonably view it as an early
(yet still incomplete) version of an AGI system.''

For managers, three things about \emph{Sparks} matter more than the
conclusion itself.

\paragraph{The honest caveats inside the paper.} The authors note that
they worked with an \emph{early, unaligned} version of GPT-4 that had
not yet been restricted by post-training safety work. They also note
that they could not inspect the training data, so they could not rule
out contamination (i.e., the possibility that an apparently novel test
question had in fact been seen in some form during training). They
selected examples they found interesting --- explicit about the
cherry-picking, but cherry-picking nonetheless. Their methodology is
deliberately qualitative.

\paragraph{The subsequent critiques.} The Schaeffer--Miranda--Koyejo
paper \citep{schaeffer2023mirage}, which won a NeurIPS 2023 Outstanding
Paper Award, showed that many ``emergent'' capabilities disappear when
you choose a smoother metric. That critique does not demolish
\emph{Sparks} --- most of the qualitative probes do not depend on a
thresholded metric --- but it does warn that the \emph{shape} of
capability growth is partly a measurement artefact.

\paragraph{What \emph{Sparks} actually established.} Even under the
most sceptical reading, the paper establishes something real: a single
system, trained only on next-token prediction on text, could produce
coherent and useful outputs across mathematics, coding, vision
(through TikZ and descriptions), theory of mind, and cross-domain
analogy, at a level that was not achievable by any previous system.
That is a qualitative shift, not a marginal improvement.

\subsection{AlphaGeometry: what a specialist system can do}

Ten months later, in January 2024, a paper in \emph{Nature} showed the
other face of the same revolution: not a generalist with emerging
abilities, but a \emph{specialist} that decisively exceeds human
experts in a narrow but intellectually demanding domain.
\citet{trinh2024alphageometry} present AlphaGeometry, a neuro-symbolic
system for solving International Mathematical Olympiad (IMO)-level
Euclidean geometry problems.

\paragraph{The architecture, in plain language.} Olympiad geometry
problems typically require a creative insight --- for example,
introducing a previously unmentioned point, line, or circle --- after
which a symbolic deduction engine can grind out the rest. AlphaGeometry
pairs a neural language model (trained to \emph{propose} such
constructions) with a fast symbolic deduction engine (which
\emph{verifies} and extends them). The neural component suggests; the
symbolic component proves. Crucially, the verifier provides a hard
ground-truth signal, so ``wrong'' proposals are filtered out
automatically.

\paragraph{The training data trick.} The most provocative aspect of
the paper, for managers thinking about the future of training data, is
\emph{how the neural component was trained}. Human geometry proofs are
scarce and hard to formalise. The authors instead generated
\emph{synthetically}: they sampled 100 million random geometric
diagrams, ran the symbolic engine to discover what was true about each
one, and turned the resulting (theorem, proof) pairs into training
data. The model was therefore trained on \emph{zero} human-written
proofs.

\paragraph{The headline result.} On 30 recent IMO geometry problems,
AlphaGeometry solved 25. The average human IMO gold medallist solves
25.9 of the corresponding problems; a silver medallist, 22.9; a bronze
medallist, 15.2. AlphaGeometry therefore sits at the gold-medal
threshold on this particular problem class.

\paragraph{The follow-up.} In July 2024, DeepMind's AlphaProof and
AlphaGeometry-2 reached a silver-medal performance on the full IMO
(across algebra, combinatorics, number theory, and geometry), getting
28 out of a possible 42 points, one below the gold-medal threshold.

\subsection{The pattern}

Put \emph{Sparks} and AlphaGeometry side by side and the pattern
becomes visible:

\begin{keyinsight}
\textbf{Whenever a domain has a fast, cheap, automatic verifier ---
mathematics, code, games with clear win conditions, formal
specifications --- systems can be trained against that verifier and
reach world-expert performance, often without any human demonstrations.
Whenever a domain lacks such a verifier --- persuasion, taste,
contested judgement, open-ended research --- progress is slower and
harder to measure.}
\end{keyinsight}

The managerial question that follows is uncomfortable and useful:
\emph{which parts of your business have (or could have) a verifier?}
Regulatory filings do. Unit tests do. Invoice reconciliation does.
Board-level strategic judgement does not. The closer a workflow is to
having a verifier, the sooner it will be automatable to expert level.

% =============================================================================
\section{From Chatbots to Agents: Measuring Autonomy}\label{sec:act3}
% =============================================================================

\subsection{What we mean by ``agent''}

A chatbot answers a single question. An \emph{agent} takes a goal,
decomposes it into sub-tasks, acts in an environment (files, web,
shell, other APIs), observes the result, and iterates. The
architectural lineage, if you want to drop the right names, starts
with \emph{ReAct} \citep{yao2023react}, which interleaved reasoning
traces with tool calls; continues with \emph{Reflexion}
\citep{shinn2023reflexion}, which added verbal self-critique; and now
includes vendor-deployed ``computer-use'' systems that drive a
browser or a desktop the way a human would, with screenshots as input
and mouse/keyboard events as output.

For a manager, the important shift is not the architecture. It is
that the unit of work has changed. The last generation of AI produced
\emph{one-shot outputs}: a translation, a classification, a
completion. The current generation produces \emph{trajectories}:
sequences of actions towards a goal, taking minutes, then hours,
potentially soon days. This is the transition from \emph{tools} to
\emph{workers}, and the speed of that transition is measurable.

\subsection{The METR task-horizon methodology}

The organisation \emph{Model Evaluation and Threat Research} (METR),
which has been running independent evaluations of frontier models
since 2023, developed a measurement that is becoming the industry's
reference \citep{metr2024horizon}.

The methodology is simple enough to explain on a whiteboard. Take a
portfolio of tasks from software engineering, research, and data
analysis. Measure how long a \emph{skilled human professional} needs
to complete each one. Then give each frontier model a chance to
complete the same tasks end-to-end, and ask: \emph{at what human-time
horizon does this model succeed at least 50\,\% of the time?} That
number --- the \emph{50\,\% task horizon} --- has a nice property:
it grows as the model gets more capable, but it does not saturate in
the uninformative way that accuracy-on-a-fixed-benchmark does.

The striking finding is this. Plotted on a log scale against release
date, the 50\,\% task horizon of frontier models has been \emph{doubling
roughly every seven months}. In 2022, the reference frontier system
could reliably complete human-minutes of work end-to-end. By 2024, it
could reliably complete human-tens-of-minutes. By 2026, the horizon is
in the low hours. If the trend continues at the same doubling time, a
task horizon of one working day arrives around 2027, and one working
week around 2029.

\subsection{How to read an extrapolation like this}

This is the moment in the lecture where managerial instincts should
fire in both directions at once.

\paragraph{The trend is real and matters.} The measurement is done by
an independent third party, on tasks that are not the tasks the models
were trained for, with human baselines collected rigorously. It is
not a vendor benchmark. The doubling time has been stable enough that
it survives the introduction of reasoning models and agentic
scaffolding.

\paragraph{Extrapolation is a claim, not a measurement.} Every
exponential trend in technology eventually bends. Hardware speed-ups
slowed as Dennard scaling ended; semiconductor feature sizes did not
continue to shrink on Moore's original cadence. The METR curve could
bend for several reasons: compute cost, data exhaustion, engineering
complexity of long-horizon reliability, or a fundamental capability
ceiling we do not yet understand. Conversely, it could accelerate
--- reasoning models, better tool use, and post-training improvements
have all produced step-changes in the past eighteen months.

\paragraph{What the plot is actually evidence for.} It is strong
evidence that (a) agent capability is moving on a measurable trend,
(b) we are not yet at any natural ceiling, and (c) the leading
indicator you should watch is the trend itself, updated quarterly,
not any single benchmark score. What it is \emph{not} evidence for is
a specific date.

\begin{keyinsight}
Managers should treat the METR curve the way they treat a leading
economic indicator --- informative about direction and magnitude, not
a forecast with confidence intervals. Plan for the doubling to
continue for roughly another eighteen months; re-plan if it bends.
\end{keyinsight}

\subsection{Adjacent evidence}

Two other benchmarks reinforce the picture. \emph{SWE-bench}
\citep{jimenez2024swebench} measures the ability of language models to
resolve real GitHub issues from open-source repositories by editing
multiple files --- a task that, two years ago, was far out of reach.
Verified variants that eliminate ambiguous issues are now routinely
solved by frontier systems above 50\,\%. \emph{GAIA}
\citep{mialon2024gaia} targets general assistance tasks that require
multimodal input, tool use, and reasoning; humans score about 92\,\%,
and frontier systems have moved from under 15\,\% at release to above
half by 2026. Both benchmarks are, by construction, more contamination-
resistant than multiple-choice tests, and both show the same trend as
the METR horizon: from unusable, through useful-with-supervision, to
outcompeting the median human junior practitioner on the narrow
task.

% =============================================================================
\section{When Will AGI Arrive?}\label{sec:act4}
% =============================================================================

\subsection{The largest expert survey ever conducted}

\citet{grace2024thousands} report the results of a January 2023 survey
of 2,778 researchers who had published in top AI venues (NeurIPS and
ICML). The survey is, to date, the largest of its kind. It asks, in
carefully designed ways, when various specific capabilities will
arrive, when ``High-Level Machine Intelligence'' (HLMI) --- defined as
machines capable of every task better and more cheaply than humans ---
will arrive, and what the consequences might be.

Three numbers from that survey are worth remembering.

\paragraph{50\,\% by 2047.} The aggregated forecast gives a 50\,\%
probability that HLMI exists by 2047. The 10\,\% probability date is
2027; the 90\,\% is after 2100. The spread is enormous. Reasonable
researchers disagree by fifty years.

\paragraph{A 13-year shift in one year.} The same survey, run in 2022,
produced a 50\,\% date of 2060. In one year, the aggregate moved
thirteen years earlier. No methodology change of that magnitude
occurred; the respondents updated.

\paragraph{38\,\% give $\geq$10\,\% to extreme outcomes.} When asked
about the probability that advanced AI leads to outcomes ``as bad as
human extinction or the permanent disempowerment of humanity,'' 38\,\%
of respondents gave a probability of at least 10\,\%. A further
substantial fraction gave probabilities above 1\,\%. This is a
minority position, but it is a large minority, and the respondents are
the people who build these systems.

\subsection{How to read an expert forecast}

Three points for the manager in the room.

\paragraph{The shift is the signal.} A 13-year jump in one year is not
a cleaner forecast replacing a noisier one. It is \emph{the same
researchers updating hard on new evidence.} When professionals
collectively revise their probabilities that much that fast, the
correct Bayesian response is not to fixate on the new point estimate
but to widen your own uncertainty and shorten your review cycle.

\paragraph{Aggregates hide the disagreement, and the disagreement is
the news.} ``Fifty percent by 2047'' reads like a forecast. It is not.
It is a weighted average over a community in which a serious minority
expects HLMI before 2030 and another serious minority expects it after
2100. Acting only on the mean is like setting your thermostat to the
average of ``too hot'' and ``too cold'' and expecting to be
comfortable.

\paragraph{The 10\,\% tail is a decision-theoretic event.} A 10\,\%
probability of an extremely bad outcome is the kind of number that,
applied to a pharmaceutical trial or a bridge design standard, would
stop the programme. If your first instinct on hearing ``only 10\,\%''
is to discount it, notice that your instinct is domain-specific. The
analogous manager in drug development would not discount it. The
question worth sitting with is whether the domain analogy should hold.

\subsection{The epistemically honest picture}

\begin{keyinsight}
We do not know when AGI arrives. We know that the people best
positioned to guess disagree by fifty years, and that the community
mean has moved thirteen years earlier in one year. The responsible
posture is neither ``this is hype; ignore it'' nor ``we have ten years;
panic'' but: \emph{build organisational capacity to update quickly when
the next strong signal arrives.}
\end{keyinsight}

% =============================================================================
\section{Managing the Transition}\label{sec:act5}
% =============================================================================

\subsection{The \emph{Science} consensus statement}

In May 2024, \emph{Science} published \emph{Managing extreme AI risks
amid rapid progress} \citep{bengio2024managing}. The signatories
include Yoshua Bengio and Geoffrey Hinton (both Turing Award laureates,
Hinton also a 2024 Nobel laureate in Physics), Stuart Russell, Andrew
Yao, Dawn Song, Pieter Abbeel, Trevor Darrell, Yuval Noah Harari,
Daniel Kahneman, and several dozen others. The paper is short, written
for a policy audience, and explicit about what the authors are and are
not claiming.

They are not claiming that AGI is imminent on a specific timeline.
They are claiming that (a) capabilities are advancing fast enough that
waiting for clear evidence of risk is not a safe default; (b) the
current research ecosystem allocates dramatically more effort to
capability than to safety; and (c) governance is lagging the
technology in ways that previous regulated industries did not tolerate.

\subsection{The four governance asks, plainly stated}

The paper makes four asks. They are worth stating in ordinary
language, because they are the substance most likely to show up in
rules and procurement standards within the next 24 months.

\begin{enumerate}[leftmargin=1.4em]
  \item \textbf{Reorient R\&D towards safety.} Devote a serious
        fraction (the paper suggests roughly a third) of frontier AI
        research effort to safety and ethical use. Today the fraction
        is small.
  \item \textbf{Adaptive governance for frontier systems.} Require
        pre-deployment capability evaluations, incident reporting, and
        the ability to pause training or deployment in response to
        risk findings.
  \item \textbf{Commensurate risk-management practices.} Hold
        developers of frontier systems to standards similar to those
        applied in aviation, pharmaceuticals, and nuclear --- including
        third-party audits and liability for harms.
  \item \textbf{International coordination.} Avoid a race to the
        bottom by aligning standards across major jurisdictions,
        including mechanisms for verification of compliance.
\end{enumerate}

Your view on whether these asks are necessary, sufficient, or
politically realistic is not the point. The point is that these are
now the reference asks that regulators are converging on. The EU AI
Act, the UK AI Safety Institute's evaluation programme
\citep{aisiuk2024}, the Bletchley and Seoul declarations, and US
executive action are all legible through this framework.

\subsection{A manager's playbook}

The honest short list --- the Monday-morning checklist --- is five
items.

\begin{enumerate}[leftmargin=1.4em]
  \item \textbf{Adopt a frontier model in one or two scoped workflows.}
        Pilot with clear success metrics and a named internal owner.
        Do not pilot AI as a general-purpose initiative; pilot it as
        a named productivity intervention on a named workflow.
  \item \textbf{Build an internal evaluation function.} A small team
        (initially one person) whose job is to run periodic capability
        checks against \emph{your} workflows, not academic benchmarks.
        This is the single highest-leverage investment in this
        decade.
  \item \textbf{Track three external signals quarterly.} The METR
        task-horizon update; the leading third-party agent benchmarks
        (SWE-bench Verified, GAIA); and the UK AISI / US AISI public
        findings. These together give you a calibrated, vendor-
        independent view.
  \item \textbf{Include AI risk in your standard enterprise risk
        review cycle.} Treat it the way you would a major regulatory
        change or a supply-chain dependency --- with an owner, a
        register, and a named escalation path.
  \item \textbf{Avoid irreversible commitments.} Multi-year exclusive
        contracts, workflows that cannot be rolled back to a
        human-in-the-loop version, and vendor lock-in on proprietary
        evaluation data are all failures of optionality in a period of
        rapid change.
\end{enumerate}

\subsection{What the honest executive posture looks like}

Neither ``ignore it, this is hype'' nor ``panic, this is existential''
is the responsible posture. Both are abdications of managerial
judgement. The responsible posture is:

\begin{itemize}[leftmargin=1.4em]
  \item \emph{Take out appropriate insurance.} You do not need to
        believe a bad outcome is likely to justify building the
        capacity to detect and respond to one.
  \item \emph{Invest in evaluation, not prediction.} You cannot time
        the arrival of AGI. You can position your organisation to
        notice when specific workflows flip from ``AI-assisted'' to
        ``AI-autonomous.''
  \item \emph{Keep AGI trajectory on the executive agenda.} Do not
        delegate this entirely to IT or procurement. It is a
        strategic question that will, within this decade, reshape
        competitive dynamics in most knowledge-work industries.
\end{itemize}

% =============================================================================
\section*{Closing: The One Thing To Remember}
\addcontentsline{toc}{section}{Closing}
% =============================================================================

If you take away nothing else from this session, take away this:

\begin{quote}
\itshape The question is not whether AGI is coming. The question is
whether your organisation has the evaluation capability to tell ---
month by month, workflow by workflow --- what these systems can and
cannot do for you. The companies that build that evaluation capability
in the next two years will be the ones that manage the transition
well. The companies that outsource the judgement to their vendors, or
their regulators, or their consultants, will not.
\end{quote}

The evidence is mixed, the timelines are contested, and the people who
build these systems disagree with each other by decades. But on one
point they agree: the rate of change has accelerated, and the window
for building institutional capacity to respond is now.

% =============================================================================
\bibliographystyle{plainnat}
\bibliography{section5}
% =============================================================================

\end{document}
